\chapter{Metodologia}

Este capítulo apresenta os passos metodológicos referentes ao trabalho. Será utilizada uma abordagem empírica, exploratória e mista, combinando métodos quantitativos e qualitativos. A metodologia será estruturada de modo a garantir rigor científico e transparência nos procedimentos adotados, bem como possibilitar a replicação do estudo \cite{sommerville2011,runeson2012}.

\section{Caracterização da Pesquisa}

Do ponto de vista de sua finalidade, esta pesquisa será classificada como aplicada, pois investigará a utilização de um método específico (USARP) em contextos reais e educacionais, visando produzir conhecimento com potencial impacto prático. Quanto aos objetivos, o estudo possuirá caráter exploratório e descritivo, uma vez que buscará compreender padrões de uso, percepção e efetividade do método, bem como descrever empiricamente os resultados observados \cite{sommerville2011,runeson2012}.

No que se refere à abordagem metodológica, o trabalho adotará uma estratégia mista, combinando análise quantitativa, baseada em estatística descritiva e Análise de Componentes Principais (PCA), e análise qualitativa, fundamentada na interpretação de evidências empíricas reportadas nos estudos primários \cite{runeson2012}.

\section{Síntese Empírica dos Estudos sobre o Método USARP}

O estudo será conduzido como uma meta-análise empírica de braço único, na qual evidências quantitativas e qualitativas provenientes de diferentes estudos empíricos sobre o método USARP serão sintetizadas de forma integrada. Diferentemente de uma meta-análise estatística comparativa clássica, este trabalho não realizará comparação entre grupos controle e experimental, mas buscará identificar padrões recorrentes, tendências temporais e estruturas latentes associadas à aplicação do método \cite{kitchenham2009,rodrigues2010}.

Os dados analisados serão extraídos de estudos empíricos conduzidos em contextos acadêmicos, envolvendo estudantes de Engenharia de Software, e contextos industriais, incluindo equipes distribuídas e ambientes de desenvolvimento remoto \cite{marques2022}.

\section{Instrumento de Coleta de Dados}

A principal fonte de dados quantitativos deste estudo consistirá em questionários estruturados aplicados após a utilização do método USARP. Esses questionários serão empregados em diversos contextos, acadêmicos e industriais, e terão como objetivo capturar a percepção dos participantes sobre a utilidade, eficácia, produtividade e aplicabilidade prática do método.

Os questionários serão organizados em dois blocos principais:

\begin{enumerate}
    \item \textbf{Bloco 1 – Percepção Geral do Método USARP}: Este bloco avaliará quatro dimensões centrais associadas ao uso do método:
    \begin{itemize}[leftmargin=1.5cm,label=\textbullet]
        \item Produtividade;
        \item Eficiência;
        \item Completude;
        \item Utilidade.
    \end{itemize}
    Para cada uma dessas dimensões, os participantes responderão em uma escala do tipo Likert de 5 pontos, variando de “Discordo totalmente” (1) a “Concordo totalmente” (5).

    \item \textbf{Bloco 2 – Percepção Específica sobre o Conteúdo das Cartas do USARP}: Este bloco será subdividido conforme os três tipos de cartas do método:
    \begin{itemize}[leftmargin=1.5cm,label=\textbullet]
        \item \textbf{Cartas Amarelas (Mecanismos de Usabilidade)};
        \item \textbf{Cartas Azuis (Requisitos)};
        \item \textbf{Cartas Laranjas (Prototipação)}.
    \end{itemize}
    Para cada carta, os participantes avaliarão diferentes aspectos como: clareza da descrição do mecanismo, contextualização, apoio à elicitação de requisitos e suporte à prototipação. As respostas serão coletadas por meio de uma escala ordinal de 6 pontos, variando de “Prejudicou” (1) a “Ajudou muito” (6).
\end{enumerate}

\section{Escala de Medida e Variáveis}

As questões dos questionários utilizarão uma escala ordinal do tipo Likert, com cinco ou seis níveis de resposta. Para viabilizar a análise estatística, as respostas textuais serão mapeadas para valores numéricos inteiros, prática comum em análises exploratórias com dados perceptivos \cite{hair2019}.

A partir dos questionários, serão definidas as seguintes variáveis analíticas: produtividade, eficiência, completude, utilidade, descrição do mecanismo de usabilidade, contexto do mecanismo de usabilidade, questão a ser explorada sobre requisitos, guia de especificação do requisito, questão sobre prototipação e guia de especificação da prototipação. Cada observação corresponderá a um grupo definido por ano e contexto (acadêmico ou industrial).

\section{Consolidação e Pré-processamento dos Dados}

Como os estudos primários apresentarão respostas agregadas por frequência, as variáveis serão calculadas por meio de médias ponderadas, considerando a distribuição das respostas em cada item do questionário, procedimento amplamente adotado em análises estatísticas exploratórias \cite{hair2019}. 

Antes da aplicação das técnicas multivariadas, os dados passarão por um processo de padronização (\textit{z-score}), garantindo que todas as variáveis contribuam igualmente para a análise, independentemente de suas escalas originais \cite{hair2019}.

\section{Análise Quantitativa Tradicional}

Inicialmente, os dados serão analisados por meio de estatística descritiva, incluindo médias, frequências, distribuições percentuais e representações gráficas. Esse tipo de análise é amplamente adotado em estudos empíricos em Engenharia de Software, conforme diretrizes metodológicas consolidadas na literatura \cite{kitchenham2009,runeson2012}, e constituirá a base analítica de grande parte dos estudos primários considerados nesta pesquisa.

A utilização de análises quantitativas tradicionais também será recorrente em revisões narrativas da área. Um exemplo representativo é o trabalho de Glass, Vessey e Ramesh (2002), que examinou 369 artigos publicados em seis periódicos relevantes de Engenharia de Software entre 1995 e 1999. Os autores classificaram manualmente os estudos com base em cinco dimensões: tópicos abordados, abordagens e métodos de pesquisa, disciplinas de referência e nível de análise. Embora o estudo tenha seguido um processo classificatório consistente, ele não foi guiado por um protocolo explícito de revisão sistemática com critérios transparentes de inclusão, exclusão ou avaliação de qualidade. A seleção dos periódicos e dos artigos (um a cada cinco) foi feita com base na percepção de relevância e conveniência, sem um protocolo pré-registrado. A síntese dos dados foi predominantemente descritiva, utilizando frequências e categorias desenvolvidas pelos próprios autores. Embora forneça uma visão ampla e útil sobre o campo de pesquisa em Engenharia de Software, a ausência de critérios formais limita a replicabilidade do estudo e o expõe a vieses de seleção e interpretação.

Em contraposição, as diretrizes para \textit{Systematic Literature Reviews} (SLR) enfatizam a necessidade de protocolos formalizados, documentação rigorosa do processo de busca e critérios explícitos de seleção e avaliação da qualidade dos estudos, como forma de mitigar vieses e aumentar a auditabilidade e a confiabilidade dos resultados \cite{kitchenham2007}. Dessa forma, ao discutir a análise quantitativa tradicional presente nos estudos primários, este trabalho utilizará o estudo de Glass et al. (2002) como referência ilustrativa de revisão narrativa/tradicional, contrastando explicitamente seus procedimentos com as práticas sistemáticas recomendadas na literatura metodológica de SLR.

Entretanto, apesar de sua ampla utilização, a análise descritiva tradicional apresentará limitações quando aplicada a conjuntos de dados compostos por múltiplas variáveis correlacionadas, dificultando a identificação de padrões globais e de relações latentes entre os atributos avaliados. Essa limitação motivará a adoção de técnicas analíticas complementares, capazes de oferecer uma visão mais integrada da estrutura dos dados.

\section{Análise de Componentes Principais}

Como técnica complementar à análise estatística tradicional, será aplicada a Análise de Componentes Principais (PCA) sobre a matriz de dados padronizada. A PCA é uma técnica estatística multivariada amplamente utilizada para redução de dimensionalidade e identificação de estruturas latentes em conjuntos de dados complexos \cite{jolliffe2002,hair2019}.

A seleção do número de componentes considerará a variância explicada individual e acumulada, a análise do \textit{scree plot} e a interpretabilidade conceitual dos componentes, conforme recomendações da literatura estatística \cite{jolliffe2002,hair2019}. Serão retidos os três primeiros componentes principais, responsáveis por explicar mais de 95\% da variância total dos dados.

\section{Análise Qualitativa e Triangulação}

Os resultados quantitativos serão interpretados à luz das evidências qualitativas reportadas nos estudos primários, incluindo análises baseadas em Grounded Theory e relatos descritivos dos participantes. Essa triangulação metodológica contribuirá para fortalecer a validade interna e externa dos achados e ampliar a compreensão dos padrões identificados \cite{runeson2012}.

\section{Ambiente Computacional}

As análises quantitativas serão realizadas utilizando a linguagem Python, com apoio das bibliotecas Pandas, NumPy, Scikit-learn e Matplotlib, garantindo reprodutibilidade, transparência metodológica e facilidade de replicação dos procedimentos adotados, conforme boas práticas em pesquisa empírica computacional \cite{kitchenham2009}.

\chapter{CONSIDERAÇÕES FINAIS}

Espera-se que a aplicação do método USARP, conforme delineado nesta pesquisa, contribua para uma elicitação de requisitos de usabilidade mais estruturada, produtiva e eficaz. A organização conceitual do método em cartas temáticas, voltadas a mecanismos, requisitos e prototipação, deverá fornecer aos participantes uma base sistemática de reflexão, favorecendo a identificação de requisitos relevantes de forma mais completa e orientada por boas práticas de design centrado no usuário. Espera-se também que a avaliação empírica permita evidenciar quais mecanismos são mais bem compreendidos e aplicados, e quais ainda apresentam barreiras de uso ou subutilização em contextos reais.

Adicionalmente, a combinação de métodos estatísticos tradicionais com Análise de Componentes Principais (PCA) poderá revelar dimensões latentes associadas ao uso do USARP, contribuindo para refinar tanto o instrumento quanto o próprio método. A comparação entre contextos acadêmicos e industriais deverá enriquecer a discussão sobre sua adaptabilidade e aplicabilidade prática. Com isso, o estudo buscará oferecer não apenas um panorama da aceitação do método, mas também insumos para sua evolução contínua, gerando conhecimento útil à comunidade de Engenharia de Requisitos e à prática profissional em projetos orientados à experiência do usuário.

\section{Cronograma Geral}

O cronograma de execução deste trabalho está dividido em dois semestres letivos. O primeiro semestre corresponde à disciplina TCC~1 (setembro de 2025 a janeiro de 2026), no qual foram desenvolvidos os capítulos iniciais do trabalho, com destaque para a fundamentação teórica e a metodologia. O segundo semestre corresponde à disciplina TCC~2 (março a julho de 2026), período em que serão realizadas as análises quantitativas e qualitativas, interpretação dos resultados e elaboração das seções finais do trabalho.

\captionsetup{width=15cm}
\captionof{table}{Cronograma de execução das atividades do TCC~1 (2025.2)}
\label{tab:cronograma_tcc1}
\IBGEtab{}{
\begin{tabular}{p{8.2cm}ccccc}
\toprule
\textbf{Atividades} & \textbf{Set/25} & \textbf{Out/25} & \textbf{Nov/25} & \textbf{Dez/25} & \textbf{Jan/26} \\
\midrule \midrule
Revisão da literatura e seleção de estudos primários & X & X &  &  &  \\
Definição dos objetivos e estruturação metodológica &  & X & X &  &  \\
Elaboração dos instrumentos de coleta &  &  & X & X &  \\
Escrita da seção Introdução & X & X &  &  &  \\
Escrita da Fundamentação Teórica &  & X & X &  &  \\
Escrita dos Trabalhos Relacionados &  &  & X & X &  \\
Escrita da Metodologia &  &  &  & X & X \\
Revisão e ajustes finais do texto do TCC~1 &  &  &  &  & X \\
\textbf{Encerramento e entrega do TCC~1} &  &  &  &  & \textbf{X} \\
\bottomrule
\end{tabular}
}{
\Fonte{Elaborado pelo autor (2025).}
}

\vspace{0.5cm}

Com a conclusão das atividades previstas no TCC~1, resta agora executar a etapa empírica do trabalho. Isso inclui a consolidação dos dados coletados, a análise estatística e qualitativa, bem como a redação dos capítulos de resultados, discussão e considerações finais. O cronograma do TCC~2 detalha essas etapas e organiza sua execução ao longo do primeiro semestre de 2026.

\vspace{0.5cm}

\captionsetup{width=15cm}
\captionof{table}{Cronograma de execução das atividades do TCC~2 (2026.1)}
\label{tab:cronograma_tcc2}
\IBGEtab{}{
\begin{tabular}{p{8.2cm}ccccc}
\toprule
\textbf{Atividades} & \textbf{Mar/26} & \textbf{Abr/26} & \textbf{Mai/26} & \textbf{Jun/26} & \textbf{Jul/26} \\
\midrule \midrule
Consolidação e padronização dos dados empíricos & X &  &  &  &  \\
Análise quantitativa tradicional (estatística descritiva) & X & X &  &  &  \\
Aplicação da Análise de Componentes Principais (PCA) &  & X & X &  &  \\
Análise qualitativa e triangulação dos dados &  &  & X & X &  \\
Interpretação dos resultados e discussão &  &  & X & X &  \\
Redação das seções de Resultados e Discussão &  &  &  & X &  \\
Redação das Considerações Finais &  &  &  & X &  \\
Revisão geral e ajustes finais &  &  &  & X & X \\
Preparação da apresentação e defesa &  &  &  &  & X \\
\textbf{Entrega e Defesa do TCC~2} &  &  &  &  & \textbf{X} \\
\bottomrule
\end{tabular}
}{
\Fonte{Elaborado pelo autor (2026).}
}